\chapter{Discussion and Conclusion}
% \label{chap:discussion}

\section{Chapter Overview}
This chapter draws together the empirical and methodological contributions of the dissertation. The central aim of the project was to evaluate automatic speech recognition (ASR) robustness under overlapping English speech, while separating the effects of overlap, background noise, model architecture and metric choice. This problem is important because realistic multi-speaker speech is neither clean nor sequential: natural conversation contains interruption, short responses, simultaneous turns and environmental noise \citep{Conversation_in_small_groups, Shriberg2001ObservationsOO, cetin06_interspeech}. Established corpora such as AMI and CHiME-6 preserve these conditions, but they do not allow the independent manipulation of overlap ratio, signal-to-noise ratio (SNR), noise type, duration and speaker count \citep{ami-corpus, carletta2005ami, chime6, watanabe2020chime6challengetacklingmultispeakerspeech}. Conversely, synthetic resources such as WHAM!, WHAMR!, LibriCSS and LibriMix have made important progress in noise, reverberation and partial overlap modelling, but they are primarily benchmark datasets rather than configurable experimental frameworks for factor-level analysis \citep{wichern2019whamextendingspeechseparation, maciejewski2020whamrnoisyreverberantsinglechannel, chen2020continuousspeechseparationdataset, cosentino2020librimixopensourcedatasetgeneralizable}.

The dissertation therefore made four connected contributions. First, it developed a reusable synthetic mixture generation framework in which overlap ratio, SNR, noise type, duration and speaker count can be controlled explicitly. Secondly, it used this framework to compare four publicly available ASR systems under matched conditions: \texttt{faster-whisper}, \texttt{WhisperX}, \texttt{wav2vec-2.0} and \texttt{Parakeet}. Thirdly, it compared synthetic trends against a real CHiME-6 evaluation subset to examine transferability. Finally, it proposed and evaluated Double-Stream Serialisation WER (DSS-WER), a beam-search approximation designed for the specific case in which a two-stream reference is compared with a single-stream hypothesis. In this way, the dissertation contributes both an experimental testbed and a metric-oriented analysis of a structural mismatch that is common in practical ASR evaluation.

% \begin{table}[htbp]
% \centering
% \caption{Summary of dissertation findings by research question.}
% \label{tab:discussion_rq_summary}
% \small
% \begin{tabularx}{\linewidth}{p{2.2cm}X}
% \toprule
% \textbf{Question} & \textbf{Conclusion} \\
% \midrule
% RQ1 & WER increases as overlap ratio rises and as SNR falls. The fitted \texttt{wav2vec-2.0} surface indicates that overlap is the dominant degradation factor, while the negative interaction term suggests that the marginal effect of one degradation source becomes smaller when the other is already severe. \\
% RQ2 & \texttt{faster-whisper} gives the lowest WER on the matched synthetic subset, while \texttt{wav2vec-2.0} gives the best latency--accuracy trade-off. Model ranking therefore depends on whether accuracy, speed or deployment constraint is prioritised. \\
% RQ3 & Synthetic trends transfer only partially to real CHiME-6 audio. The real-data ranking changes substantially, showing that controlled synthetic benchmarking is diagnostically useful but insufficient as a sole proxy for real conversational robustness. \\
% RQ4 & DSS-WER is close to ordinary WER in non-overlapping clips and significantly lower overall in overlapping clips. However, its behaviour is model-dependent, and it should be interpreted as a task-specific complementary metric rather than a universal replacement for WER. \\
% \bottomrule
% \end{tabularx}
% \end{table}

\section{Discussion of Results}
\subsection{Overlap and Noise as Degradation Factors}
The results in \Cref{chap:statistical_analysis} confirm the central assumption of the dissertation: overlap and noise should be treated as first-class experimental variables rather than incidental properties of a test set. In the full \texttt{wav2vec-2.0} sweep, WER increased with both higher overlap ratio and lower SNR. This is consistent with previous work showing that multi-speaker speech violates the single-active-speaker assumption on which many conventional ASR evaluations implicitly depend \citep{galibert2013methodologies, chen2020continuousspeechseparationdataset, von_Neumann_2023}. It also agrees with the broader observation that realistic listening conditions are shaped by the joint effect of competing talkers and acoustic background rather than either factor alone \citep{Conversation_in_small_groups, wu2018characteristics}.

The fitted surface in \Cref{chap:statistical_analysis} estimated
\begin{equation}
    \widehat{\mathrm{WER}} =
    0.029 + 1.001\,\mathrm{OVR} + 0.274\,z - 0.328\,(\mathrm{OVR}\times z),
\end{equation}
where \(z\) is the linear amplitude noise ratio derived from SNR. The positive overlap and noise terms show that both factors are harmful. The larger coefficient for overlap suggests that, in the tested \texttt{wav2vec-2.0} conditions, overlap contributes more strongly to WER than additive noise. The negative interaction is also important. It suggests an antagonistic rather than purely additive relationship: once the acoustic scene is already difficult because of strong overlap, the additional measured penalty of noise becomes smaller, and conversely once the noise is severe, the additional measured penalty of overlap is less separable. This does not mean that noise becomes harmless. Rather, it means that the two error sources partly saturate the same recognition capacity.

This interpretation is consistent with the motivation behind LibriCSS and more recent realistic overlap simulation work. \citet{chen2020continuousspeechseparationdataset} show that continuous, partially overlapping speech requires different evaluation assumptions from fully separated or fully overlapped mixtures. \citet{yang2022simulatingrealisticspeechoverlaps} further argue that random overlap is insufficient because real overlap has conversational structure. The present results support both positions: overlap is not merely another form of noise, and its effect depends on the structure of the test condition. The public noise category also produced the highest WER for \texttt{wav2vec-2.0}, which suggests that irregular, speech-like or highly variable backgrounds may be more disruptive than more stationary transport noise. This agrees with the dissertation's use of DEMAND as a practical source of environment-specific noise rather than a single undifferentiated background condition \citep{demand}.

\subsection{Model Behaviour Under Controlled Conditions}
The matched synthetic comparison shows that \texttt{faster-whisper} achieved the lowest mean WER, followed by \texttt{WhisperX}, \texttt{wav2vec-2.0} and \texttt{Parakeet}. This ordering broadly reflects the advantages of large-scale weakly supervised encoder--decoder ASR under diverse audio conditions, as represented by Whisper-style systems \citep{radford2022robustspeechrecognitionlargescale, github_openai_faster-whisper, bain2026whisperx_repo}. At the same time, the timing benchmark showed that \texttt{wav2vec-2.0} was substantially faster than the other models while still producing a moderate WER, which makes it attractive when throughput or real-time processing is a constraint \citep{baevski2020wav2vec20frameworkselfsupervised, facebook_wav2vec2-large-960h}.

The model comparison should not be read as a universal leaderboard. It is instead a statement about the specific operating points tested in this dissertation. \texttt{faster-whisper} and \texttt{WhisperX} had the best synthetic WER, but both also showed severe hallucination-style outliers in some clips. These failures were qualitatively different from the worst \texttt{wav2vec-2.0} cases, which tended to resemble deletion or near-silence under severe noise rather than fluent repeated text. This distinction matters because WER alone collapses qualitatively different failure modes into the same scalar penalty \citep{morris04_interspeech, sasindran2023hevalnewhybridevaluation, ellis2026werunawareassessingasr, hughes2025problemwithwer}. A model that hallucinates fluent text and a model that omits speech may have similar WER in a particular clip, but they create different risks for downstream users.

The \texttt{Parakeet} result is especially instructive. On the matched synthetic subset it had the highest WER, yet in the real CHiME-6 subset it became the best model among the four. Since Parakeet is a conformer/transducer-style model designed for efficient ASR, its behaviour may reflect a different balance between acoustic robustness, decoding assumptions and segmentation behaviour \citep{sekoyan2025canary1bv2parakeettdt06bv3efficient, nvidia/parakeet-tdt-0.6b-v3}. The dissertation does not provide enough evidence to claim why the reversal occurs, but it does show why model selection should not rely on a single synthetic benchmark. The practical conclusion is that synthetic testing is valuable for controlled diagnosis, while deployment-oriented selection still requires evaluation on audio that resembles the intended use case.

\subsection{Synthetic-to-Real Transferability}
The transferability analysis is one of the most important findings of the dissertation. The synthetic benchmark successfully exposed the effects of overlap and SNR under controlled conditions, but the real CHiME-6 evaluation changed the model ranking. This confirms a familiar trade-off in ASR evaluation. Synthetic data provides control and repeatability; real corpora provide ecological validity \citep{ami-corpus, carletta2005ami, chime6, watanabe2020chime6challengetacklingmultispeakerspeech}. Neither replaces the other.

Several factors explain the gap. The synthetic mixtures used in this dissertation were generated from clean read speech sources, including LibriSpeech and VCTK, with additive DEMAND noise \citep{librispeech, voxceleb, demand}. Although this design supports clean factorial analysis, it does not fully reproduce far-field recording, room reverberation, microphone placement, spontaneous turn-taking, laughter, hesitations, interruption, speaker movement or the social organisation of a real conversation. CHiME-6, by contrast, contains unsegmented distant-microphone recordings in a natural domestic setting \citep{watanabe2020chime6challengetacklingmultispeakerspeech}. The higher real-data WER for most models therefore does not invalidate the synthetic benchmark; it identifies the boundary of what the benchmark controls.

This result aligns with the direction of recent CHiME challenges and meeting recognition research. The field has increasingly moved from clean recognition towards distant, multi-speaker, multi-array and generalisable meeting transcription \citep{BARKER2017605, chime7, chime8_task1, cornell2024chime8dasrchallengegeneralizable}. It has also increasingly combined separation, diarisation and recognition in modular or end-to-end pipelines \citep{raj2020integrationspeechseparationdiarization, chen2020continuousspeechseparationconformer, vonneumann2024meetingrecognitioncontinuousspeech}. The present dissertation complements that work by showing what can be learned from a controlled benchmark, while also demonstrating why controlled benchmarks must be validated against real audio before they are used for deployment claims.

\subsection{Metric Choice and the Role of DSS-WER}
The DSS-WER analysis addresses a narrower but important methodological problem. Standard WER assumes that the reference and hypothesis can be compared as single word sequences. In overlapping speech, however, the reference may naturally contain two simultaneous speaker streams, while many ASR systems still output a single serialised transcript. Existing multi-speaker metrics such as cpWER, ORC-WER, MIMO-WER and DI-cpWER formalise several important speaker-attributed and speaker-agnostic settings, but the double-stream-reference to single-stream-hypothesis case remains awkward in practice \citep{galibert2013methodologies, von_Neumann_2023, vonneumann2024meetevaltoolkitcomputationword, von_Neumann_2025, Word_Error_Rate_Definitions}.

DSS-WER was introduced to address this case by allowing any interleaving of the two reference streams that preserves the within-speaker word order. The empirical results support the motivation, but with qualifications. In non-overlapping clips, DSS-WER remained very close to ordinary WER, which is desirable because there is no interleaving to reward. In overlapping clips, DSS-WER was significantly lower than WER overall, especially for \texttt{wav2vec-2.0}. This suggests that ordinary WER can over-penalise a hypothesis that serialises overlapping speech in a valid but different order.

However, DSS-WER is not a universal improvement. It is an approximate metric, and for \texttt{Parakeet} it produced a higher mean score than WER in the overlapping condition. This implies that the beam search can be sensitive to hypothesis quality and search ambiguity. The metric should therefore be used as an additional diagnostic rather than as a replacement for established measures. Its strongest use case is not to declare one model universally better, but to test whether a model's apparent WER penalty is partly caused by reference serialisation rather than lexical error. This is consistent with recent metric literature, which emphasises that evaluation measures must be matched to the output structure and application being studied \citep{von_Neumann_2023, vonneumann2024meetevaltoolkitcomputationword, von_Neumann_2025}.

\section{Main Contributions}
The first contribution is the construction of a controlled synthetic mixture generation framework. By varying overlap ratio, SNR, noise category, duration and speaker count explicitly, the framework makes it possible to ask factor-level questions that cannot be answered cleanly from a single natural corpus. This sits between fixed synthetic benchmarks and real meeting corpora: it preserves enough control for causal interpretation while remaining connected to realistic ASR degradation factors \citep{chen2020continuousspeechseparationdataset, cosentino2020librimixopensourcedatasetgeneralizable, yang2022simulatingrealisticspeechoverlaps}.

The second contribution is the empirical comparison of heterogeneous ASR systems under identical audio and reference conditions. The dissertation shows that accuracy, latency and failure mode do not align perfectly. \texttt{faster-whisper} performed best by synthetic WER, \texttt{wav2vec-2.0} gave the best speed, and \texttt{Parakeet} transferred best to the real CHiME-6 subset. This provides a more nuanced view than a single aggregate score.

The third contribution is the synthetic-to-real comparison. The dissertation shows that controlled synthetic mixtures are useful for understanding the effect of specific degradation factors, but that they only partially predict model ranking on real conversational audio. This finding is methodologically important because it prevents the synthetic framework from being over-claimed. The framework is best understood as a diagnostic instrument rather than a replacement for real-world evaluation.

The fourth contribution is the proposal and empirical evaluation of DSS-WER. While related to the broader family of multi-speaker WER definitions, DSS-WER targets a particular structural mismatch that arises when double-stream references meet single-stream ASR hypotheses. Its results show that interleaving-aware scoring can change the interpretation of overlapped clips, while also revealing the need for further work on exactness, efficiency and robustness.

\section{Limitations}
Several limitations constrain the interpretation of the results. First, the synthetic speech material is based on clean read-speech corpora rather than spontaneous conversation. LibriSpeech and VCTK provide clean transcripts and speaker diversity, but they do not fully capture conversational disfluency, interruptions, backchannels or turn-taking behaviour \citep{librispeech, voxceleb, Shriberg2001ObservationsOO}. This means the synthetic benchmark is well suited for controlled acoustic analysis but less suited for modelling all linguistic properties of meetings.

Secondly, the acoustic simulation is deliberately simple. Additive DEMAND noise supports controlled SNR and noise-category analysis, but reverberation, microphone geometry, speaker movement and spatial separation are not modelled in the main experiment \citep{demand, maciejewski2020whamrnoisyreverberantsinglechannel}. This limits the direct comparability of synthetic results to far-field datasets such as CHiME-6.

Thirdly, the main synthetic experiments fix the maximum number of concurrent speakers at two and use approximately 60-second clips. This makes the DSS-WER formulation and the factorial analysis tractable, but real meetings can involve more speakers, longer contexts and complex turn-taking patterns \citep{ami-corpus, chime6, cornell2024chime8dasrchallengegeneralizable}. The conclusions should therefore be read as strongest for two-speaker overlap rather than for all multi-party ASR.

Fourthly, the real-world transferability evaluation uses 100 clips from a single CHiME-6 session. This is sufficient to show that transfer is imperfect, but not sufficient to estimate general real-world performance across different rooms, speakers, microphones and social settings. A broader real-data evaluation would be needed before making robust deployment claims.

Fifthly, model coverage is limited to four public ASR systems and not every model was evaluated on the full 4000-clip synthetic grid. The matched 700-clip subset supports fair cross-model comparison, but it does not provide a full factorial model-by-condition analysis for every system. The \texttt{wav2vec-2.0} full-grid results are therefore useful for understanding condition effects, while the matched subset is more appropriate for comparing models.

Finally, DSS-WER is currently a beam-search approximation. Its behaviour in the non-overlapping and overlapping tests is encouraging, but the \texttt{Parakeet} case shows that approximation error or search ambiguity can affect the score. The metric also focuses on lexical interleaving and does not directly evaluate speaker attribution, timing accuracy or diarisation quality. It should therefore be considered alongside WER, ORC-WER and other multi-speaker metrics rather than in isolation \citep{von_Neumann_2023, vonneumann2024meetevaltoolkitcomputationword, von_Neumann_2025}.

\section{Future Research Directions}
\subsection{More Realistic Synthetic Mixture Generation}
The most immediate future direction is to improve the realism of the synthetic generator while preserving factor-level control. This could include room impulse responses, reverberation, spatialised sources, microphone-array simulation and more varied background conditions, following the motivation behind WHAMR!, CHiME-style distant ASR and meeting recognition benchmarks \citep{maciejewski2020whamrnoisyreverberantsinglechannel, BARKER2017605, watanabe2020chime6challengetacklingmultispeakerspeech}. A second extension would be to replace purely random utterance scheduling with learned or corpus-derived overlap patterns. Recent work on realistic overlap simulation suggests that modelling the temporal structure of interruptions and conversational turns can improve multi-talker ASR training and evaluation \citep{yang2022simulatingrealisticspeechoverlaps}.

\subsection{Broader Real-World Validation}
The synthetic-to-real gap should be tested across more corpora and sessions. AMI, additional CHiME-6 sessions, CHiME-7 and CHiME-8 style datasets would allow the framework to be evaluated against a wider distribution of rooms, microphones, speaker counts and social contexts \citep{carletta2005ami, chime7, cornell2024chime8dasrchallengegeneralizable}. This would make it possible to distinguish which synthetic trends are stable across real settings and which are artefacts of the generation process.

\subsection{Extension to Speaker-Attributed and Multi-Output ASR}
Future work should extend the evaluation beyond single-stream hypotheses. Modern meeting transcription increasingly combines speech separation, diarisation and recognition, or performs speaker-attributed ASR directly \citep{raj2020integrationspeechseparationdiarization, chen2020continuousspeechseparationconformer, kanda2021endtoendspeakerattributedasrtransformer, kanda2022transcribetodiarizeneuralspeakerdiarization, cui2024improvingspeakerassignmentspeakerattributed, raj2024speakerattributionsurt, zheng2025dncasrendtoendtrainingspeakerattributed}. Evaluating these systems would require a richer metric suite, including speaker-aware WER, diarisation-aware scores and timing-sensitive measures. This would also allow the synthetic framework to support research into whether errors arise from separation, diarisation, speaker attribution or lexical recognition.

\subsection{Further Development of DSS-WER}
DSS-WER should be developed in three directions. First, it should be compared against exact solutions on short examples to quantify approximation error and tune the beam-search hyperparameters. Secondly, the search should be stress-tested under low-quality hypotheses, since the present results suggest that the approximation is less reliable when the model output is poor. Thirdly, DSS-WER should be integrated with existing meeting-evaluation toolkits so that its behaviour can be compared systematically with ORC-WER, cpWER, MIMO-WER and DI-cpWER \citep{vonneumann2024meetevaltoolkitcomputationword, Word_Error_Rate_Definitions}. This would help establish whether DSS-WER is best understood as a diagnostic tool, a reporting metric or an intermediate component inside a larger scoring framework.

\subsection{Training and Adaptation Experiments}
The framework developed here was used primarily for evaluation, but it could also be used for training or fine-tuning. A natural next step is to train models on controlled mixtures and test whether robustness improves specifically in the degradation conditions represented during training. Prior work suggests that realistic overlap simulation and integrated separation-recognition pipelines can improve multi-talker performance \citep{yang2022simulatingrealisticspeechoverlaps, raj2020integrationspeechseparationdiarization, chen2020continuousspeechseparationconformer}. The present framework would allow such gains to be analysed by factor: for example, whether training at high overlap improves low-overlap recognition, or whether exposure to one noise category transfers to another.

\section{Final Conclusion}
This dissertation has shown that robust ASR evaluation under overlapping English speech requires both experimental control and realistic validation. The controlled synthetic benchmark demonstrated that overlap ratio and SNR have clear, interpretable effects on WER, with overlap emerging as the stronger degradation factor in the full \texttt{wav2vec-2.0} analysis. The model comparison showed that no single system dominated all criteria: \texttt{faster-whisper} achieved the lowest synthetic WER, \texttt{wav2vec-2.0} achieved the best inference efficiency, and \texttt{Parakeet} performed best on the limited real CHiME-6 subset. The real-data analysis further showed that synthetic rankings cannot be assumed to transfer directly to natural conversation.

The dissertation also showed that metric choice is part of the scientific problem. Ordinary WER remains useful, but it can misrepresent overlapping speech when the reference is naturally multi-stream and the hypothesis is single-stream. DSS-WER provides a principled first step towards evaluating this mismatch by allowing valid interleavings of two reference streams. Its empirical behaviour supports the motivation, while its limitations show that future work is needed before it can be treated as a mature scoring standard.

Overall, the main conclusion is that overlapping speech should not be treated as a single benchmark condition. It is a structured degradation factor whose effect depends on noise, model architecture, decoding behaviour and evaluation metric. A controlled synthetic framework, when paired with real-corpus validation and careful metric selection, offers a clearer and more honest way to evaluate ASR robustness in the conversational settings where ASR is increasingly expected to work.


